\documentclass[12pt, twoside]{article}
%\documentclass[12pt, twoside]{article}
\usepackage[letterpaper, margin=1in, headsep=0.2in]{geometry}
\setlength{\headheight}{0.6in}
%\usepackage[english]{babel}
\usepackage[utf8]{inputenc}
\usepackage{microtype}
\usepackage{amsmath}
\usepackage{amssymb}
%\usepackage{amsfonts}
\usepackage[nomessages]{fp} %\FPeval{\var-name}{2*sin(pi/6)}
\usepackage{siunitx} %units in math. eg 20\milli\meter
\usepackage{yhmath} % for arcs, overparenth command
\usepackage{tikz} %graphics
\usetikzlibrary{quotes, angles, arrows, arrows.meta}
\usepackage{pgfplots}
\usepgfplotslibrary{fillbetween}
\pgfplotsset{compat=1.16}
\usepackage{graphicx} %consider setting \graphicspath{{images/}}
\usepackage{parskip} %no paragraph indent
\usepackage{enumitem}
\usepackage{multicol}
\usepackage{venndiagram}

\usepackage{fancyhdr}
\pagestyle{fancy}
\fancyhf{}
\renewcommand{\headrulewidth}{0pt} % disable the underline of the header
\raggedbottom
\hfuzz=2mm %suppresses overfull box warnings

\usepackage{hyperref}

\fancyhead[LE]{\thepage}
\fancyhead[RO]{Markscheme}
\fancyhead[LO]{La Scuola d'Italia / Dr.\ Huson / IB Math 11 \\* 21 October 2026}

\begin{document}
\subsubsection*{2.3 Pre-Quiz: Sequences, Lines, Radicals and Exponents --- Markscheme}

\begin{enumerate}[itemsep=0.3cm]

%--- Problem 1: Geometric sequence 24, 12, 6: r, u_6, why S_inf exists; find r from u_1 and S_inf [6 marks] ---
\item[1.] [Maximum mark: 6]

\textbf{(a)} $r = \frac{12}{24} = \frac{1}{2}$ \hfill \textbf{A1}

\textbf{(b)} $u_{6} = 24\left(\frac{1}{2}\right)^{5}$ \hfill \textbf{M1}

$u_{6} = \frac{24}{32} = \frac{3}{4}$ \hfill \textbf{A1}

\emph{Accept $0.75$. FT from (a).}

\textbf{(c)} $|r| < 1$ \quad (since $r = \frac{1}{2}$) \hfill \textbf{R1}

\emph{Accept ``$-1 < r < 1$'' or ``$r$ is between $-1$ and $1$''. Do not
accept ``the terms get smaller'' alone.}

\textbf{(d)} $\frac{10}{1 - r} = 15$ \hfill \textbf{M1}

$1 - r = \frac{10}{15} = \frac{2}{3}$

$r = \frac{1}{3}$ \hfill \textbf{A1}

\newpage

%--- Problem 2: Points P, Q; midpoint, gradient of PQ, perpendicular bisector L, point (k, 9) on L [9 marks] ---
\item[2.] [Maximum mark: 9]

\textbf{(a)} $M = \left(\frac{-4 + 6}{2},\ \frac{1 + (-3)}{2}\right)$

$M = (1,\ -1)$ \hfill \textbf{A1A1}

\emph{A1 for each coordinate.}

\textbf{(b)} $m_{PQ} = \frac{-3 - 1}{6 - (-4)}$ \hfill \textbf{M1}

$m_{PQ} = -\frac{4}{10} = -\frac{2}{5}$ \hfill \textbf{A1}

\textbf{(c)} $m_{L} = \frac{5}{2}$ \hfill \textbf{A1}

\emph{FT from (b): the negative reciprocal of their gradient.}

\textbf{(d)} $y - (-1) = \frac{5}{2}(x - 1)$ \hfill \textbf{M1}

$y + 1 = \frac{5}{2}(x - 1)$ \hfill \textbf{A1}

\emph{M1 for substituting their gradient and their $M$ into any form of the
equation of a line. Accept $y = \frac{5}{2}x - \frac{7}{2}$ or $5x - 2y = 7$.
FT from (a) and (c).}

\textbf{(e)} $9 + 1 = \frac{5}{2}(k - 1)$ \hfill \textbf{M1}

$k - 1 = 4$

$k = 5$ \hfill \textbf{A1}

\emph{M1 for substituting $(k, 9)$ into their equation of $L$. FT from (d).}

\newpage

%--- Problem 3: f(x) = 2x - 3; intercepts; parallel g through origin; perpendicular h through (2, 6); intersection of f and h; x-intercept of h [12 marks] ---
\item[3.] [Maximum mark: 12]

\textbf{(a)} $y$-intercept $-3$ \hfill \textbf{A1}

$2x - 3 = 0$, so $x = \frac{3}{2}$ \hfill \textbf{A1}

\emph{Accept $(0, -3)$ and $\left(\frac{3}{2}, 0\right)$.}

\textbf{(b)} gradient of $g$ is $2$ \hfill \textbf{A1}

$g(x) = 2x$ \hfill \textbf{A1}

\emph{The gradient may be seen in the expression.}

\textbf{(c)} gradient of $h$ is $-\frac{1}{2}$ \hfill \textbf{A1}

$y - 6 = -\frac{1}{2}(x - 2)$ \hfill \textbf{M1}

$h(x) = -\frac{1}{2}x + 7$ \hfill \textbf{A1}

\emph{M1 for substituting their gradient and $(2, 6)$ into any form of the
equation of a line. Accept $h(x) = -\frac{1}{2}(x - 2) + 6$.}

\textbf{(d)} $2x - 3 = -\frac{1}{2}x + 7$ \hfill \textbf{M1}

$\frac{5}{2}x = 10$

$x = 4$ \hfill \textbf{A1}

$y = 2(4) - 3 = 5$, so the point is $(4,\ 5)$ \hfill \textbf{A1}

\emph{FT from (c).}

\textbf{(e)} $-\frac{1}{2}x + 7 = 0$ \hfill \textbf{M1}

$R = (14,\ 0)$ \hfill \textbf{A1}

\emph{Accept $x = 14$. FT from (c).}

\newpage

%--- Problem 4: 4^x * 2^(x+1) = 2^(ax+b); solve = 16; solve 9^x = 27; show (sqrt12 + sqrt3)^2 = 27; 6/(3 - sqrt3) in the form a + b sqrt3 [13 marks] ---
\item[4.] [Maximum mark: 13]

\textbf{(a)} $4^{x} = \left(2^{2}\right)^{x} = 2^{2x}$ \hfill \textbf{M1}

$2^{2x} \times 2^{x + 1} = 2^{3x + 1}$

$a = 3$, \quad $b = 1$ \hfill \textbf{A1A1}

\textbf{(b)} $2^{3x + 1} = 2^{4}$ \hfill \textbf{M1}

$3x + 1 = 4$

$x = 1$ \hfill \textbf{A1}

\emph{M1 for writing $16$ as $2^{4}$ and equating exponents. FT from (a).}

\textbf{(c)} $\left(3^{2}\right)^{x} = 3^{3}$ \hfill \textbf{M1}

$2x = 3$ \hfill \textbf{A1}

$x = \frac{3}{2}$ \hfill \textbf{A1}

\emph{M1 for writing both sides as powers of $3$. Accept $x = 1.5$.}

\textbf{(d)} $\sqrt{12} = 2\sqrt{3}$ \hfill \textbf{M1}

$\left(2\sqrt{3} + \sqrt{3}\right)^{2} = \left(3\sqrt{3}\right)^{2}$

$= 9 \times 3$ \hfill \textbf{A1}

$= 27$ \hfill \textbf{AG}

\emph{Accept expansion: $12 + 2\sqrt{36} + 3$ (M1), $12 + 12 + 3$ (A1).}

\textbf{(e)} $\frac{6}{3 - \sqrt{3}} \times \frac{3 + \sqrt{3}}{3 + \sqrt{3}}$ \hfill \textbf{M1}

$= \frac{6\left(3 + \sqrt{3}\right)}{9 - 3}$ \hfill \textbf{A1}

$= 3 + \sqrt{3}$ \hfill \textbf{A1}

\emph{A1 for the denominator $6$. Accept $a = 3$, $b = 1$.}

\end{enumerate}
\end{document}
